\label{sec:heterogeneity}

\subsection{State-Level Variation}

The national aggregate (20 additional competitive districts) conceals
substantial state-level heterogeneity.  Some states gain many competitive
districts under algorithmic maps; others gain few or even lose competitive
districts.  Understanding this variation reveals which features of a state's
geography and demography drive the competitiveness difference.

Table~\ref{tab:state-gains} presents the states with the largest gains in
competitive districts under algorithmic maps.

\begin{table}[h]
\centering
\caption{States with Largest Gains in Competitive Districts (Algorithmic vs.\ Enacted)}
\label{tab:state-gains}
\begin{tabular}{lcccc}
\toprule
State & Districts & Algorithmic & Enacted & Gain \\
\midrule
North Carolina & 14 & 5 & 2 & +3 \\
Texas & 38 & 10 & 6 & +4 \\
Ohio & 15 & 5 & 2 & +3 \\
Florida & 28 & 8 & 5 & +3 \\
Georgia & 14 & 4 & 2 & +2 \\
Virginia & 11 & 4 & 2 & +2 \\
Michigan & 13 & 4 & 2 & +2 \\
Wisconsin & 8 & 3 & 1 & +2 \\
Pennsylvania & 17 & 5 & 3 & +2 \\
Arizona & 9 & 3 & 2 & +1 \\
\bottomrule
\end{tabular}
\end{table}

\subsection{Characteristics of High-Gain States}

States that gain the most competitive districts under algorithmic maps share
several characteristics.

\textbf{Geographic diversity.}  High-gain states tend to have a mix of
urban cores, suburban rings, and rural areas---the geographic configuration
that enables mapmakers to create safe seats by separating these communities.
In North Carolina, for example, enacted maps pack Democratic-leaning urban
areas of the Triangle and Charlotte into a small number of safe Democratic
seats, leaving the surrounding suburban areas in safe Republican seats.
Algorithmic maps that cross urban-suburban boundaries naturally produce
more competitive districts.

\textbf{Enacted map quality.}  States with heavily gerrymandered enacted
maps (in either direction) show the largest gains.  Texas and Ohio, which
received Republican-drawn maps widely criticized as partisan gerrymanders,
show gains of 4 and 3 competitive seats respectively.  Illinois, which
received a Democratic gerrymander, shows a gain of 2 competitive seats.

\textbf{Partisan balance.}  States with near-even partisan balance in the
presidential vote tend to show larger gains because the underlying electorate
supports more competitive districts regardless of map design.  The algorithmic
map captures this latent competitiveness; enacted maps may suppress it.

\subsection{States with Few or No Gains}

Several states show no gain or a small loss of competitive districts under
algorithmic maps.  These tend to be:
\begin{itemize}
  \item \textbf{Highly partisan states:} States where the electorate is
    overwhelmingly Democratic (Massachusetts, Maryland) or Republican
    (Wyoming, West Virginia) leave little room for competitive districts
    regardless of the map.
  \item \textbf{Small-delegation states:} States with one or two congressional
    districts rarely have the geographic range to produce a competitive
    district; the state partisan balance determines the outcome almost
    entirely.
  \item \textbf{States where enacted maps are already compact:} A handful of
    states where enacted maps use independent commissions or court-ordered
    compactness requirements show little difference from algorithmic output.
    California under its CRC map is the clearest example.
\end{itemize}

\subsection{The Role of Geographic Sorting}

The geographic sorting of Democrats into urban cores and Republicans into
suburban and rural areas~\citep{chen2013unintentional,rodden2019why} is the
single largest factor limiting the number of competitive algorithmic districts.
Even a perfectly compact map must draw district lines through the urban-suburban
boundary, and that boundary coincides with sharp partisan gradients.  The
algorithm produces competitive districts when it must cross that boundary to
achieve population balance, but in states where population is already
distributed in ways that allow purely urban or purely rural districts to achieve
balance, competitive districts are geometrically constrained.

This suggests that the gain in competitive districts from algorithmic
redistricting is bounded: the algorithm can eliminate artificial packing and
cracking but cannot overcome the underlying geographic separation of partisan
voters.  The 31\% increase we document represents the recoverable competitiveness
that gerrymandering suppresses, not the theoretical maximum.
